\begin{proposition}[Refuted: the diameter condition cannot be dropped for generic connected graphs]
\label{C:connected-convex-warning}
For a generic pair of induced graph embeddings with connected source
and target, strong convexity of both source images, injectivity on
components, and bounded nontrivial point paths do not guarantee
convexity of the old copies in the repaired beta target. This failure
can occur at $n=3$, $s=45$, $k=11$, even though these numbers satisfy
the numerical C.1 condition for $a_0=2$, $a_1=8/5$.
The graph configuration below is not asserted to arise among the
arrows of the prescribed C.1 construction.
\end{proposition}

\begin{proof}
More generally fix an odd $s=4k+1$ with $k\geq2$, and an integer
$L\geq s$. Let $S$ be the path with vertices $0,\ldots,L$, let
$T$ be the cycle with vertices $0,\ldots,2L+1$ in cyclic order, and
define
\[
 f(r)=r,\qquad g(r)=L+1+r\qquad(0\leq r\leq L).
\]
Both image paths have length $L$, strictly less than half the cycle
length $L+1$. Every shortest path in $T$ between two vertices of
either image is therefore the unique path in that image. Thus both
embeddings are strongly convex. Source and target are connected, so
their component maps are injective. The point-transition relation is
the disjoint union of the two-vertex paths
$r\mapsto L+1+r$. In particular any common bound $n\geq2$ is valid.
The two actual point maps are also increasing for the natural orders
on the displayed vertex sets.

Use the actual-path-size beta construction. Its vertex set is
$\{0,\ldots,L\}\times\mathbb Z/s$, with copy $i$ given by
\[
 \alpha_i f(r)=(r,u_i),\qquad
 \alpha_i g(r)=(r,v_i),\qquad
 (u_i,v_i)=(\beta_i^2(0),\beta_i^2(1)).
\]
The graph has a vertical path of length $L$ in every column, and for
every index pair $\{u_i,v_i\}$ it has the two end bridges
\[
 (L,u_i)(0,v_i),\qquad (0,u_i)(L,v_i).
\]
The unordered index pairs are the edges of the cycle
\[
 0,1,\ldots,2k,4k,4k-1,\ldots,2k+1,0.
\]
Because $s$ is odd, starting at $(0,0)$ and following only the end
bridges along this whole index cycle ends at $(L,0)$ after $s$
edges. The two vertices belong to copy zero, where their intrinsic
distance is $L$. This bridge path leaves copy zero. If $L>s$, it
is shorter than the old distance; if $L=s$, it has the old distance
and is an outside shortest path. To verify the latter assertion,
any path that uses an interior vertex of a vertical column and
connects distinct column ends must traverse a full column and has
length at least $L$. The bridge subgraph on the two ends of each
column is the connected double cover of an odd $s$-cycle, hence a
$2s$-cycle; the two ends of column zero are opposite and have bridge
distance $s$. Thus the new distance is exactly $\min\{L,s\}$.

For a strict shortening at numerically admissible parameters, take
$L=46$, $n=3$, $k=11$, and $s=45$. The source is a $46$-edge path,
the old target a $94$-cycle, and the repaired target has $2115$
vertices and $2160$ edges. The indicated old distance $46$ becomes
$45$. The standalone program
\texttt{akr\_connected\_convex\_counterexample.py} verifies both
old convex embeddings, the point-path bound, all neighboring copy
identities, the explicit $45$-edge bridge path, and its exact new
distance. Its JSON certificate specifies every map and every vertex
of the shorter path.

The numerical compatibility is not a counterexample to universal
C.1. In particular, if the displayed old cycle were one of the
protected first-stage circles, either proper half-cycle arrow image
would contradict Proposition~\ref{C:protected-nonadjacent}. The
specific realization of such a configuration among later generated
objects remains an additional requirement.
\end{proof}
